% Part: second-order-logic
% Chapter: syntax-and-semantics
% Section: terms-formulas

\documentclass[../../../include/open-logic-section]{subfiles}

\begin{document}

\olfileid{sol}{syn}{frm}

\olsection{Term lan \printtoken{P}{formula}}

Kaya ing logika orde kapisan, ekspresi logika orde kapindho
dibangun saka kosakata dhasar kang ngemot \emph{!!{variable}s},
\emph{!!{constant}s}, \emph{!!{predicate}s}, lan kadhang uga
\emph{!!{function}s}. Saka unsur-unsur iku, bebarengan karo
konektif logis, kuantor, lan simbol tandha wacan kayata
kurung lan koma, dibentuk \emph{term} lan \emph{!!{formula}s}.
Bedane, saliyane variabel kanggo objek, logika orde kapindho
uga ngemot variabel kanggo relasi lan fungsi, lan ngidini
kuantifikasi tumrape. Mula simbol logis logika orde kapindho
yaiku simbol logika orde kapisan, ditambah:

\begin{enumerate}
\item Himpunan kang !!{denumerable}s saka !!{variable}s relasi
  orde kapindho kanggo saben aritas~$n$: $\Obj V_0^n$,
  $\Obj V_1^n$, $\Obj V_2^n$, \dots
\item Himpunan kang !!{denumerable}s saka !!{variable}s fungsi
  orde kapindho:
  $\Obj u_0^n$, $\Obj u_1^n$, $\Obj u_2^n$, \dots
\end{enumerate}

Kaya dene kita nggunakake $x$, $y$, $z$ minangka metavariabel
kanggo variabel orde kapisan $\Obj v_i$, kita bakal nggunakake
$X$, $Y$, $Z$, lan sateruse minangka metavariabel kanggo
$\Obj V_i^n$, lan $u$, $v$, lan sateruse minangka metavariabel
kanggo~$\Obj u_i^n$.

\begin{explain}
Simbol nonlogis ing basa orde kapindho ditemtokake kanthi cara
kang padha kaya ing basa orde kapisan: kanthi ndhaptar
!!{constant}s, !!{function}s, lan !!{predicate}s.

Ing logika orde kapisan, !!{identity}~$\eq$ biasane dilebokake.
Ing logika orde kapisan, simbol nonlogis saka basa~$\Lang{L}$
wigati supaya kita bisa nyatakake perkara kang narik kawigaten.
Mesthi ana !!{sentence}s kang ora nggunakake simbol nonlogis,
nanging yen mung nganggo~$\eq$, angel nyatakake perkara kang
narik kawigaten. Ing logika orde kapindho, amarga kita nduweni
variabel relasi lan fungsi tanpa wates, kita bisa nyatakake
apa wae kang bisa dinyatakake ing basa orde kapisan sanajan ora
nduweni simbol nonlogis kang disedhiyakake kanthi mligi.
\end{explain}

\begin{defn}[Term Orde Kapindho]
Himpunan \emph{term orde kapindho} saka~$\Lang L$,
$\TrmSOL[L]$, ditetepake kanthi nambahake klausa ing ngisor
iki marang \olref[fol][syn][frm]{defn:terms}:
\begin{enumerate}
\item Yen $u$ iku variabel fungsi $n$-papan lan $t_1$, \dots,
  $t_n$ iku term, mula $\Atom{u}{t_1, \ldots, t_n}$ iku term.
\end{enumerate}
\end{defn}

\begin{explain}
Dadi, term orde kapindho wujude kaya term orde kapisan, nanging
ing papan kang ngemot !!a{function}~$\Obj{f^n_i}$ ing term orde
kapisan, term orde kapindho bisa ngemot variabel fungsi~$\Obj{u^n_i}$.
\end{explain}

\begin{defn}[\usetoken{s}{formula} Orde Kapindho]
Himpunan \emph{!!{formula}s orde kapindho}~$\FrmSOL[L]$
saka basa~$\Lang L$ ditetepake kanthi nambahake klausa-klausa
ing ngisor iki marang \olref[fol][syn][frm]{defn:formulas}:
\begin{enumerate}
\item Yen $X$ iku variabel predikat $n$-papan lan $t_1$, \dots,
  $t_n$ iku term orde kapindho saka~$\Lang L$, mula
  $\Atom{X}{t_1,\ldots, t_n}$ iku !!{formula} atomik.

\tagitem{prvAll}{Yen $!A$ iku !!a{formula} lan $u$ iku variabel fungsi,
  mula $\lforall[u][!A]$ iku !!a{formula}.}{}

\tagitem{prvAll}{Yen $!A$ iku !!a{formula} lan $X$ iku variabel predikat,
  mula $\lforall[X][!A]$ iku !!a{formula}.}{}

\tagitem{prvEx}{Yen $!A$ iku !!a{formula} lan $u$ iku variabel fungsi,
  mula $\lexists[u][!A]$ iku !!a{formula}.}{}

\tagitem{prvEx}{Yen $!A$ iku !!a{formula} lan $X$ iku variabel predikat,
  mula $\lexists[X][!A]$ iku !!a{formula}.}{}
\end{enumerate}
\end{defn}

\end{document}
